% Plain Dharma — XeLaTeX preamble for the PRINT-SHOP edition (A5, 12pt).
%
% This is the file you hand a copy shop. It differs from pdf-preamble-print.tex
% (the KDP-shaped 5×8) in four ways:
%
%   1. A5 trim (148×210mm) — the A-series twin of the old 5.25×8.25: same
%      height, 15mm wider. Two A5 pages tile an A4 sheet exactly, which is how
%      the shop prints it (2-up, duplex, cut down the middle, bind).
%   2. 12pt body instead of 10pt. The measure grows to ~114mm, so this lands at
%      ~62 characters per line — inside the 60–75 comfort band, not merely bigger.
%   3. White paper, no \pagecolor. Shops print on stock they already have; a
%      cream flood would cost a full-coverage ink charge on every page.
%   4. NO bleed and NO cover pages. The covers are a separate 2-page color file
%      (see build-printshop-pdf.ts) because the shop runs them on card in a
%      different pass — a single mixed file gets printed entirely in color by
%      mistake, which is the expensive way to find out.
%
% Margins are set in build-printshop-pdf.ts, not here: inner 20mm for the
% binding, outer/top/bottom generous enough that a 3mm trim never touches type.

% Garamond Libre — same family as every other edition.
\setmainfont{GaramondLibre}[
  Path=__FONTS_DIR__/,
  Extension=.otf,
  UprightFont=*-Regular,
  ItalicFont=*-Italic,
  BoldFont=*-Bold,
  BoldItalicFont=*-BoldItalic,
  Ligatures=TeX,
  Numbers=OldStyle
]

% Brand palette. Everything here prints as gray in the shop's B&W pass; the
% values are kept so a shop that quotes color interior gets the real colors.
\definecolor{ink}{HTML}{2A2520}
\definecolor{muted}{HTML}{6B5F50}
\definecolor{saffron}{HTML}{C7651C}

\color{ink}
\linespread{1.18}

\usepackage{titlesec}
\titleformat{\chapter}[display]
  {\normalfont\LARGE\bfseries\color{ink}}
  {}
  {0pt}
  {\LARGE}
\titlespacing*{\chapter}{0pt}{28pt}{22pt}

\titleformat{\section}
  {\normalfont\Large\bfseries\color{ink}}
  {}
  {0pt}
  {}

% Blockquote (per-chapter drops): italic, muted, indented.
\renewenvironment{quote}{%
  \begin{list}{}{%
    \leftmargin=1.2em\rightmargin=1.2em\itemindent=0pt%
  }\item\itshape\color{muted}}{%
  \end{list}}

% Illustrations. build-printshop-pdf.ts rewrites each standalone image into an
% explicit \begin{center} block, so the default \includegraphics width here is
% only a fallback — see the centering note in that script.
\usepackage{graphicx}
\setkeys{Gin}{width=0.66\linewidth,keepaspectratio}

% Pali diacritics fallback (same as every other edition).
\usepackage{newunicodechar}
\newfontfamily\palifallback{STIX Two Text}
\newunicodechar{ṇ}{{\palifallback ṇ}}
\newunicodechar{ṭ}{{\palifallback ṭ}}

% Let long source URLs break across lines — the measure is 114mm and an Access
% to Insight path is far wider. url.sty ships with BasicTeX; xurl does not.
\usepackage{url}
\Urlmuskip=0mu plus 1mu\relax
\def\UrlBreaks{\do\/\do\-\do\.\do\_\do\?\do\&\do\=\do\#\do\:}

\usepackage{enumitem}
\setlist{itemsep=0.15em,topsep=0.3em}

\setlength{\parindent}{0pt}
\setlength{\parskip}{0.6em}

\hyphenation{Mindfulness Loving-Kindness Dharma Buddha Buddhist}

% Pad the book out to a multiple of 4 pages. 2-up duplex puts four A5 pages on
% one A4 sheet, so a page count that isn't divisible by 4 leaves the shop with a
% half-used final sheet and an ambiguous cut. build-printshop-pdf.ts measures the
% first pass and substitutes the count here on the second.
__PAD_PAGES__

% Print hyperlink convention: link text renders as plain ink. Nobody can click
% a booklet.
\AtBeginDocument{%
  \hypersetup{
    colorlinks=false,
    pdfborder={0 0 0},
    pdfauthor={Gautama Buddha; translated by Claude Opus; edited by Alex Miller},
    pdfsubject={Six Foundational Buddhist Teachings in Plain Modern English},
    pdfkeywords={Buddhism, suttas, dharma, meditation, mindfulness},
    pdfcreator={pandoc + XeLaTeX (print-shop edition)}
  }%
}
